% =====================================================================================
% Paste-ready pieces for the progress report. Every number is a macro from
% results/progress/numbers.tex (regenerated by the notebook), so re-running the
% notebook and replacing numbers.tex updates the text.
%
% Preamble:
%   \usepackage{graphicx,booktabs,tikz}  \usetikzlibrary{positioning,arrows.meta}
%   \input{numbers.tex}
% Figures: copy figures/progress/*.pdf next to the .tex (they are 5.5in wide = \textwidth).
%
% Lines marked  % CHECK  depend on the outcome of the run: read the printed numbers and
% keep the sentence that matches.
% =====================================================================================


% ------------------------------- Data Processing -------------------------------------
\paragraph{Data.} We stream the January 2025 Lichess rated-games dump (CC0) and keep rated rapid games
between two humans that ended normally, whose players are within 200 points, and whose average rating
falls in one of five 200-point bands (1100--2100). To spread the sample over time, each eligible game
is kept with probability \keepProb{} until every band has \nGamesBand{} games; this scanned
\nScanned{} games played between \spanStart{} and \spanEnd{} UTC. From each game we sample up to 12
positions, skipping the first 10 plies, moves made with under 30\,s on the clock, and forced moves.
Each position is flipped to the mover's view and encoded as 17 binary $8{\times}8$ planes (12 piece,
4 castling, 1 en passant) plus 5 one-hot rating-band planes. The label is the played move as a
from--to index in $[0, 4096)$ (promotions collapsed). This gives \nTotal{} positions (about
\avgLegal{} legal moves each), split 80/10/10 \emph{by game} into \nTrain{}/\nVal{}/\nTest{}.
Only one-way hashes of usernames are kept, to measure player overlap: \seenFrac\% of test positions
come from a player who also appears in training.

\paragraph{New data for final testing.} The February 2025 dump was never used for training or tuning;
we built \nNew{} test positions from it (\nGamesNewBand{} games per band) with the same pipeline.

% Figure: statistics + one cleaned example
\begin{figure}[h]
  \centering
  \includegraphics[width=\textwidth]{fig_data.pdf}
  \caption{Cleaned data. (a) Positions per band are balanced by construction. (b) Ply of the sampled
  positions. (c) A cleaned training example, shown from the mover's side.}
  \label{fig:data}
\end{figure}

% ------------------------------- Results table ---------------------------------------
\begin{table}[h]
  \centering\small
  \caption{Move-matching on held-out January games and on February games. ResNet rows: mean
  $\pm$ s.d. over \nSeeds{} seeds; other models one seed.}
  \label{tab:results}
  \begin{tabular}{lrrrrr}
    \toprule
    Model & Params & Top-1 (\%) & Top-3 (\%) & CE (nats) & Feb.\ top-1 (\%) \\
    \midrule
    Uniform random legal move & -- & \topOneRand & -- & -- & -- \\
    Capture heuristic (no learning) & 0 & \topOneHeur & -- & -- & \newHeur \\
    Linear (from + to) & \linParams & \topOneLin & \topThreeLin & \lossLin & \newLin \\
    v1 CNN, FC head (proposal) & \vOneParams & \topOneVOne & \topThreeVOne & \lossVOne & \newVOne \\
    ResNet, no band input & \cnnParams & \topOneNc{}$\pm$\topOneNcSd & \topThreeNc & \lossNc & \newNc \\
    \textbf{ResNet, band input} & \cnnParams & \textbf{\topOneCnn}$\pm$\topOneCnnSd & \topThreeCnn & \lossCnn & \newCnn{}$\pm$\newCnnSd \\
    \bottomrule
  \end{tabular}
\end{table}

% ------------------------------- Baseline text ----------------------------------------
\paragraph{Baselines.} The capture heuristic reaches \topOneHeur\% top-1, against \topOneRand\% for a
random legal move. Its accuracy falls from \heurLow\% in the lowest band to \heurHigh\% in the highest,
which suggests that weaker players take available material more often. % CHECK the direction
The linear model scores a move as a from-square term plus a to-square term (Fig.~\ref{fig:baselines}).
It reaches \topOneLin\% and needs no tuning beyond the shared optimiser settings.

% ------------------------------- Primary model text -----------------------------------
\paragraph{Primary model.} The CNN from the proposal (v1) put 97\% of its \vOneParams{} parameters in one
$2048{\to}4096$ fully connected layer. Its validation loss bottomed out at epoch \bestEpVOne{} and then
rose while training loss kept falling. We replaced that layer with a fully convolutional policy head
and residual blocks (Fig.~\ref{fig:model}), which gives \cnnParams{} parameters. We train with Adam
(lr $10^{-3}$, cosine decay, weight decay $10^{-4}$, batch 512) for \nEpochs{} epochs and keep the
checkpoint with the lowest validation loss. On held-out games the model reaches \topOneCnn\% top-1
(95\% game-bootstrap CI \topOneCnnLo--\topOneCnnHi), and \newCnn\% on the unseen February games.

% ------------------------------- Band conditioning ------------------------------------
\paragraph{Does the rating band help?} Adding the band input changes top-1 by \bandGain{} points
(95\% CI \bandGainLo{} to \bandGainHi) and cross-entropy by \bandGainLoss{} nats (positive = better).
Shuffling the band at test time gives \topOneCnnShuf\% versus \topOneCnn\%. Giving the model the
matching band rather than a mismatched one changes top-1 by \diagAdvTop{} points
(Fig.~\ref{fig:bands}b). On average, switching the band from 1100--1300 to 1900--2100 moves the
predicted distribution by a total-variation distance of only \meanTV{} (largest \bandTV{} over \nTV{}
positions, Fig.~\ref{fig:qual}c--d).
% CHECK: un-comment ONE of the next two sentences.
% (a) if the CI on \bandGain is entirely above 0:
% The band is a small but real signal, so the conditioned model is kept as the primary model.
% (b) if the CI includes 0 or is below 0:
% So far the band adds little beyond the position itself. Most of what separates rating levels may lie in
% how often players blunder, which a single position shows weakly. Testing this is the main goal
% before the final report (see challenges).

% ------------------------------- Figures ---------------------------------------------
\begin{figure}[h]
  \centering
  \includegraphics[width=\textwidth]{fig_curves.pdf}
  \caption{(a) Cross-entropy per epoch; dots mark the kept checkpoints. The v1 FC-head network
  overfits after epoch \bestEpVOne{}. (b) Validation top-1 for all models.}
  \label{fig:curves}
\end{figure}

\begin{figure}[h]
  \centering
  \includegraphics[width=\textwidth]{fig_bands.pdf}
  \caption{(a) Test top-1 per rating band. (b) Test positions from each true band (rows) scored with
  every band input (columns); colour is the difference from the row mean, bold marks the matching band.}
  \label{fig:bands}
\end{figure}

\begin{figure}[h]
  \centering
  \includegraphics[width=\textwidth]{fig_qualitative.pdf}
  \caption{Qualitative results. (a, b) Typical correct and incorrect predictions. (c, d) The test
  position (out of \nTV{}) where the band input changes the prediction most; this one is cherry-picked,
  and the average effect is much smaller (mean TV \meanTV).}
  \label{fig:qual}
\end{figure}

% ------------------------------- Challenges ------------------------------------------
\paragraph{Challenges and plan.}
(1) \emph{Data volume:} the proposal targeted about 2M positions; streaming and parsing on Colab limited
this run to \nTotal{}. The pipeline stops on a per-band quota, so scaling up only needs more runtime.
(2) \emph{Overfitting} of the v1 head, addressed with the convolutional head above.
(3) \emph{Player overlap:} splitting by game leaves \seenFrac\% of test positions with a player seen in
training. On positions from unseen players only, top-1 is \topOneUnseen\% (\topOneSeen\% with seen
players). The final report will split by player.
(4) \emph{Narrow time window:} all games come from \spanStart{} to \spanEnd{} UTC. The final run will draw
from several months.
(5) \emph{Band signal:} next steps are per-band models (as in Maia) and adding clock time as an input.
The proposal's Maia comparison ($\sim$50\% top-1) used millions of games per band and the recent move history; we
use a single position and far less data.
